\documentclass[12pt, letterpaper]{article}

\usepackage[margin=1in]{geometry}
\usepackage{graphicx}
\usepackage{booktabs}
\usepackage{amsmath}
\usepackage{hyperref}
\usepackage{caption}
\usepackage{float}
\usepackage{microtype}
\usepackage{parskip}
\usepackage{xcolor}
\usepackage{enumitem}
\usepackage[numbers, sort&compress]{natbib}

\hypersetup{colorlinks=true, linkcolor=blue, citecolor=blue, urlcolor=blue}

\title{\textbf{CNN vs.\ Transformer for 3D Brain Tumor Segmentation:\\
       A Controlled Comparison with Boundary-Aware Loss}}
\author{Devansh Sharma\\
        CAP 5516 --- Medical Image Computing, Spring 2026}
\date{}

\begin{document}
\maketitle
\thispagestyle{empty}

% ─────────────────────────────────────────────────────────────────────────────
\begin{abstract}
Accurate segmentation of brain tumors from multi-modal MRI is a critical step in
treatment planning for glioblastoma.
This project compares a CNN-based segmentation model (DynUNet, an nnU-Net-style
network) against a Vision Transformer (SwinUNETR) on the Medical Segmentation
Decathlon Task~01 brain tumor dataset.
Both models are trained under identical conditions and evaluated on per-class
Dice Similarity Coefficient (DSC) and 95th-percentile Hausdorff Distance (HD95)
across three tumor sub-regions.
We find that DynUNet achieves slightly higher mean DSC (0.741 vs.\ 0.732) while
SwinUNETR achieves better mean HD95 (6.19\,mm vs.\ 6.94\,mm), confirming
a well-known CNN--Transformer trade-off: CNNs excel at voxel-overlap but
Transformers produce smoother, better-localised boundaries.
As our original contribution, we introduce a \emph{boundary-aware loss}
that combines standard DiceCE loss with a weighted Hausdorff Distance
Transform loss applied periodically during training.
This single modification reduces DynUNet's weakest-class HD95
(Peritumoral Edema) from 9.88\,mm to 7.10\,mm --- a 28\% improvement that
brings it below SwinUNETR's score on the same class, without adding any
model parameters or increasing inference time.
\end{abstract}

% ─────────────────────────────────────────────────────────────────────────────
\section{Introduction}

Glioblastoma multiforme (GBM) is among the most aggressive primary brain tumors,
with a median survival of roughly 15 months.
Precise volumetric segmentation from MRI guides surgical planning, radiotherapy
dose contouring, and longitudinal monitoring.
Manual delineation is time-consuming and subject to inter-rater variability,
making automatic segmentation a high-priority clinical need.

The field has been shaped by two major architectural families.
Convolutional encoder--decoder networks---exemplified by U-Net \cite{ronneberger2015} and
its 3D descendants---leverage local spatial inductive biases and have been
the workhorse of medical image segmentation for almost a decade.
Vision Transformers \cite{dosovitskiy2020} capture long-range dependencies
through self-attention, at the cost of higher memory and a weaker spatial prior.
SwinUNETR \cite{tang2022} combines hierarchical shifted-window attention with a
CNN decoder, representing the current Transformer state-of-the-art on brain
tumor benchmarks.

This project frames the comparison as a controlled experiment: identical data
splits, preprocessing, optimizers, and patch sizes, so that architectural
differences drive the metric gap rather than training tricks.
We then ask whether a lightweight loss-level intervention can close the
CNN's boundary-accuracy gap without switching to a heavier architecture.

\textbf{Original contributions.}
\begin{enumerate}[leftmargin=*, label=(\arabic*)]
    \item A fair, reproducible head-to-head comparison of DynUNet and SwinUNETR
          on the MSD Task~01 dataset using the MONAI framework.
    \item A boundary-aware loss function (DiceCE + weighted HausdorffDT,
          applied every $k=10$ training steps) that specifically targets
          HD95 without inflating training cost.
    \item Empirical analysis showing that the boundary loss closes
          DynUNet's ED HD95 gap to SwinUNETR by 28\%, while preserving
          sub-10\,ms-per-volume inference speed.
\end{enumerate}

% ─────────────────────────────────────────────────────────────────────────────
\section{Related Work}

\paragraph{nnU-Net and DynUNet.}
Isensee et al.\ \cite{isensee2021} showed that a well-tuned standard U-Net
consistently outperforms most hand-crafted architectures.
MONAI's DynUNet implements the same self-configuring design: kernel sizes and
strides are derived from the dataset's median voxel spacing, and auxiliary
deep supervision heads are attached at each decoder resolution.

\paragraph{SwinUNETR.}
Tang et al.\ \cite{tang2022} pre-trained a Swin Transformer encoder on 5{,}000
unlabelled CT/MRI volumes and fine-tuned it for the BraTS 2021 segmentation
challenge, achieving state-of-the-art Dice scores.
Its hierarchical window attention is resolution-agnostic, making it straightforward
to apply to arbitrary patch sizes.

\paragraph{Boundary-aware losses.}
Several works have noted that Dice loss is insensitive to boundary errors.
Kervadec et al.\ \cite{kervadec2019} proposed a contour-based loss term;
Karimi \& Gholipour \cite{karimi2019} derived a differentiable Hausdorff
distance approximation.
MONAI provides \texttt{HausdorffDTLoss}, which back-propagates through a
distance-transform approximation of HD.
We leverage this directly as a regularisation term.

% ─────────────────────────────────────────────────────────────────────────────
\section{Methodology}

\subsection{Problem Formulation}
Given a 4-channel MRI volume $\mathbf{X} \in \mathbb{R}^{4 \times H \times W \times D}$
(T1, T1ce, T2, FLAIR modalities), the goal is to predict a voxel-level label map
$\hat{\mathbf{Y}} \in \{0,1,2,3\}^{H \times W \times D}$, where labels correspond
to: 0 = background, 1 = necrotic core / non-enhancing tumor (NCR/NET),
2 = peritumoral edema (ED), 3 = enhancing tumor (ET).

\subsection{Dataset}
We use the \textbf{Medical Segmentation Decathlon Task~01: Brain Tumour} dataset
\cite{antonelli2022}, which contains 484 multi-modal MRI cases in 4D NIfTI
format (four modalities stacked along the channel axis).
Segmentation labels follow the convention 0/1/2/3.
Cases are split deterministically (seed = 42) into
338 training / 72 validation / 74 test cases (70/15/15\%).

\subsection{Preprocessing}
All volumes are resampled to 1.0\,mm isotropic spacing, reoriented to RAS,
and z-score normalised per modality using only non-zero voxels (to exclude
background bias).
During training, random $128^3$ patches are extracted; validation and test
use sliding-window inference with a $128^3$ window and batch size~4.

\subsection{Model Architectures}

\paragraph{DynUNet (CNN baseline).}
A 5-level encoder--decoder with residual blocks, instance normalisation, and
leaky ReLU activations.
Kernel sizes are $3\times3\times3$ throughout; encoder strides are
$[1,2,2,2,2]$.
Deep supervision is enabled: auxiliary outputs at resolutions
$\times\frac{1}{2}$ and $\times\frac{1}{4}$ contribute to training loss
with weights $1.0$ and $0.5$.
Trainable parameters: \textbf{16.5\,M}.

\paragraph{SwinUNETR (Transformer).}
A Swin Transformer encoder with window size 7, 4 stages, and feature
dimension 48, connected to a CNN decoder via skip connections.
No pre-training was used; the model was trained from scratch to maintain
a fair comparison.
Trainable parameters: \textbf{62.2\,M} (3.8$\times$ DynUNet).

\subsection{Loss Functions}

\paragraph{Baseline loss (both models).}
\begin{equation}
    \mathcal{L}_{\text{DiceCE}} = \mathcal{L}_{\text{Dice}} + \mathcal{L}_{\text{CE}}
\end{equation}
where $\mathcal{L}_{\text{Dice}}$ is the soft multi-class Dice loss and
$\mathcal{L}_{\text{CE}}$ is cross-entropy, both with softmax activations.

\paragraph{Boundary-aware loss (DynUNet ablation).}
\begin{equation}
    \mathcal{L}_{\text{boundary}} =
        \mathcal{L}_{\text{DiceCE}} + \lambda \cdot \mathcal{L}_{\text{HD}}
\end{equation}
where $\mathcal{L}_{\text{HD}}$ is MONAI's \texttt{HausdorffDTLoss}
(a differentiable distance-transform approximation of HD95) and $\lambda = 0.5$.
Computing 3D distance transforms at every step is prohibitively slow
(${\sim}16$\,s/batch vs.\ ${\sim}2$\,s for DiceCE alone).
We therefore apply $\mathcal{L}_{\text{HD}}$ only every $k=10$ training steps,
falling back to plain $\mathcal{L}_{\text{DiceCE}}$ otherwise.
This reduces the per-epoch overhead by ${\approx}5\times$ while still
exposing the model to boundary gradients regularly throughout training.
Deep supervision weights $\{1.0, 0.5, 0.25\}$ are applied to each auxiliary
head output.

\subsection{Training Protocol}
Both models use identical settings unless stated otherwise (Table~\ref{tab:training}).

\begin{table}[H]
\centering
\caption{Training hyperparameters shared by all three runs.}
\label{tab:training}
\begin{tabular}{ll}
\toprule
\textbf{Hyperparameter} & \textbf{Value} \\
\midrule
Optimizer           & Adam \\
Learning rate       & $1\times10^{-3}$ \\
Weight decay        & $1\times10^{-5}$ \\
LR schedule         & Cosine Annealing \\
Batch size          & 2 \\
Patch size          & $128^3$ \\
Max epochs          & 300 (DynUNet baseline \& SwinUNETR); 150 (DynUNet+BL) \\
Val interval        & 5 epochs \\
Cache rate          & 1.0 \\
Data augmentation   & Random crop to patch; normalisation only \\
Hardware            & NVIDIA A100 80\,GB (HPC, single GPU) \\
\bottomrule
\end{tabular}
\end{table}

Checkpoints are saved whenever validation mean Dice improves.
Final test evaluation uses the best validation checkpoint.

% ─────────────────────────────────────────────────────────────────────────────
\section{Experiments and Results}

\subsection{Quantitative Results}

Table~\ref{tab:results} reports per-class and mean DSC and HD95 on the
74-case held-out test set.
Figure~\ref{fig:bars} visualises the same metrics as grouped bar charts.

\begin{table}[H]
\centering
\caption{Test-set segmentation results. Best per column in \textbf{bold}.
         $\uparrow$ = higher is better; $\downarrow$ = lower is better.}
\label{tab:results}
\resizebox{\textwidth}{!}{%
\begin{tabular}{lcccccccc}
\toprule
 & \multicolumn{3}{c}{\textbf{Dice} $\uparrow$} & & \multicolumn{3}{c}{\textbf{HD95 (mm)} $\downarrow$} & \\
\cmidrule(lr){2-4} \cmidrule(lr){6-8}
\textbf{Model} & NCR/NET & ED & ET & Mean & NCR/NET & ED & ET & Mean \\
\midrule
DynUNet (baseline)        & \textbf{0.794} & \textbf{0.598} & \textbf{0.831} & \textbf{0.741} & 6.60 & 9.88 & \textbf{4.33} & 6.94 \\
SwinUNETR                 & 0.793 & 0.587 & 0.815 & 0.732 & 7.41 & 7.27 & 3.89 & \textbf{6.19} \\
DynUNet + BoundaryLoss    & 0.786 & 0.580 & 0.801 & 0.722 & \textbf{6.39} & \textbf{7.10} & 6.03 & 6.50 \\
\midrule
\multicolumn{9}{l}{\small $^\dagger$ Inference time: DynUNet 0.008\,s/vol; SwinUNETR 0.122\,s/vol. GPU: DynUNet 10.3\,GB; SwinUNETR 17.0\,GB.}\\
\bottomrule
\end{tabular}}
\end{table}

\begin{figure}[H]
    \centering
    \includegraphics[width=\linewidth]{figures/fig1_per_class_metrics.png}
    \caption{Per-class Dice (left) and HD95 (right) for all three models on the
             test set. The green arrow on the HD95 plot highlights the 28\%
             reduction in ED HD95 achieved by the boundary-aware loss.}
    \label{fig:bars}
\end{figure}

\begin{figure}[H]
    \centering
    \includegraphics[width=0.75\linewidth]{figures/fig4_ed_hd95_spotlight.png}
    \caption{ED (Peritumoral Edema) HD95 comparison. The boundary-aware loss
             reduces DynUNet's ED HD95 from 9.88\,mm to 7.10\,mm, dropping below
             SwinUNETR's score of 7.27\,mm on this class.}
    \label{fig:spotlight}
\end{figure}

\subsection{Computational Efficiency}

Table~\ref{tab:compute} summarises the computational footprint of each model.

\begin{table}[H]
\centering
\caption{Computational comparison on the same A100 hardware.}
\label{tab:compute}
\begin{tabular}{lcccc}
\toprule
\textbf{Model} & \textbf{Params} & \textbf{GPU Mem} & \textbf{Inference} & \textbf{Epochs} \\
\midrule
DynUNet (baseline)     & 16.5\,M & 10.3\,GB & 0.008\,s/vol & 300 \\
SwinUNETR              & 62.2\,M & 17.0\,GB & 0.122\,s/vol & 300 \\
DynUNet + BoundaryLoss & 16.5\,M & 10.3\,GB & 0.008\,s/vol & 150 \\
\bottomrule
\end{tabular}
\end{table}

SwinUNETR requires $3.8\times$ more parameters, $1.65\times$ more GPU memory,
and is $15\times$ slower at inference than DynUNet.
The boundary-aware loss adds \emph{zero} inference overhead since
$\mathcal{L}_{\text{HD}}$ is only active during training.

\subsection{Training Dynamics}

Figure~\ref{fig:curves} shows validation Dice and training loss across epochs.

\begin{figure}[H]
    \centering
    \includegraphics[width=\linewidth]{figures/fig3_training_curves.png}
    \caption{Left: validation Dice over training epochs. Right: training loss.
             The DynUNet+BoundaryLoss curve (green) shows instability early on
             due to HausdorffDT gradient spikes, but stabilises and converges
             by epoch 120. SwinUNETR (orange) converges faster and plateaus
             around epoch 50. DynUNet baseline (blue dashed) converges slowly
             but reaches the highest peak validation Dice of 0.742 by epoch 300.
             The boundary-loss training loss starts at ${\sim}11$ (DiceCE +
             HausdorffDT combined) before dropping sharply.}
    \label{fig:curves}
\end{figure}

% ─────────────────────────────────────────────────────────────────────────────
\section{Analysis and Discussion}

\paragraph{Dice vs.\ HD95 trade-off.}
DynUNet achieves the highest mean DSC (0.741) but the worst mean HD95 (6.94\,mm).
SwinUNETR flips this: lower DSC (0.732) but better HD95 (6.19\,mm).
This confirms the established finding that Transformers model global context more
effectively, leading to spatially coherent predictions with fewer stray voxels
at boundaries.
The gap on ET HD95 is especially striking: SwinUNETR scores 3.89\,mm vs.\
DynUNet's 4.33\,mm, a 10\% improvement on the most clinically relevant sub-region.

\paragraph{Why does DynUNet have strong NCR/ET Dice?}
DynUNet's local convolutions are well-suited to the compact, high-contrast
enhancing tumor and necrotic core, where strong local texture cues are sufficient.
ED is larger and more diffuse, making its boundary harder to delineate without
long-range context --- which explains why SwinUNETR holds an advantage on ED HD95.

\paragraph{Effect of the boundary-aware loss.}
Adding $\mathcal{L}_{\text{HD}}$ to DynUNet's training objective reduces ED HD95
by 2.78\,mm (28\%) and mean HD95 by 0.44\,mm, narrowing the gap to SwinUNETR
from 0.75\,mm to 0.31\,mm.
The improvement is specific to the classes where DynUNet originally struggled:
NCR HD95 improves by 0.21\,mm and ED HD95 by 2.78\,mm.
ET HD95, where DynUNet was already competitive, worsens slightly (4.33 $\to$
6.03\,mm), suggesting a trade-off: the boundary loss pushes the model toward
smoother ED boundaries at the cost of some ET boundary precision.
Mean DSC decreases modestly (0.741 $\to$ 0.722), which is the typical
DSC--HD trade-off when adding boundary terms.

\paragraph{Training instability with HausdorffDT.}
Two sharp validation Dice drops are visible at epochs 30 and 95
(Figure~\ref{fig:curves}).
These correspond to moments when the HausdorffDT term produces large gradient
spikes on cases with near-zero predicted probability for a class.
The model recovers in both cases, suggesting the Adam optimiser with cosine
annealing is robust enough to handle these perturbations.
Using $k=10$ (applying the boundary loss every 10 steps rather than every step)
was essential: a frequency of $k=1$ required ${\sim}16$\,s/iteration vs.\
${\sim}2$\,s/iteration, making 300-epoch training computationally infeasible.

\paragraph{Inference speed as a practical consideration.}
SwinUNETR is 15$\times$ slower at inference (0.122\,s vs.\ 0.008\,s per volume)
due to the quadratic cost of attention in 3D patches.
For a clinical deployment processing hundreds of MRI studies per day, DynUNet's
sub-10\,ms inference is a meaningful practical advantage.

\paragraph{Limitations.}
\begin{itemize}[leftmargin=*]
    \item \textbf{No pre-training for SwinUNETR.}
          The original SwinUNETR paper uses self-supervised pre-training on 5{,}000
          volumes, which can add 3--5\% DSC.
          Our from-scratch training likely underestimates SwinUNETR's ceiling.
    \item \textbf{DynUNet+BL trained for only 150 epochs.}
          The baseline ran 300 epochs and was still slowly improving at convergence.
          The boundary-loss model may not have fully converged, so the reported
          DSC gap between baseline and BL may shrink with additional training.
    \item \textbf{No data augmentation beyond random crop.}
          Standard augmentations (random flip, intensity jitter, elastic deformation)
          were omitted to keep the comparison controlled, but they would likely
          improve all models.
    \item \textbf{Single random seed.}
          Results are reported for one training run per model.
          Variance across seeds was not measured.
\end{itemize}

% ─────────────────────────────────────────────────────────────────────────────
\section{Conclusions}

This project presented a controlled comparison of CNN and Transformer architectures
for 3D brain tumor segmentation on the MSD Task~01 dataset.
The main findings are:

\begin{enumerate}[leftmargin=*, label=(\arabic*)]
    \item \textbf{CNNs win on Dice; Transformers win on HD95.}
          DynUNet scores 0.741 mean DSC vs.\ SwinUNETR's 0.732, but SwinUNETR
          achieves 6.19\,mm mean HD95 vs.\ DynUNet's 6.94\,mm.
          The choice of metric matters: for tumour volume estimation, DynUNet is
          preferable; for surgical margin delineation, SwinUNETR is better.

    \item \textbf{A boundary-aware loss bridges the HD95 gap.}
          Periodic application of \texttt{HausdorffDTLoss} during DynUNet training
          reduces ED HD95 by 28\% and overall mean HD95 by 6.3\%,
          closing the gap to SwinUNETR from 0.75\,mm to 0.31\,mm with no change
          to the model architecture or inference pipeline.

    \item \textbf{Efficiency favours DynUNet.}
          At $15\times$ faster inference and $3.8\times$ fewer parameters,
          DynUNet remains the more practical choice for deployment,
          particularly after boundary-loss fine-tuning.
\end{enumerate}

\textbf{Future work} could include: self-supervised pre-training for SwinUNETR,
a full ablation of boundary loss weight $\lambda$ and frequency $k$,
anisotropic distance transform weighting to prioritise clinically critical
sub-regions, and extension to the full BraTS 2021 dataset (1{,}251 cases).

% ─────────────────────────────────────────────────────────────────────────────
\section{Use of External Resources}

\paragraph{Frameworks and libraries.}
This project uses \textbf{MONAI} \cite{monai2020} (v1.3) for data loading,
transforms, model definitions (DynUNet, SwinUNETR), loss functions
(DiceCELoss, HausdorffDTLoss), and sliding-window inference.
\textbf{PyTorch} \cite{paszke2019} is used as the deep learning backend.
\textbf{TensorBoard} is used for training logging and curve visualisation.

\paragraph{Models.}
The DynUNet and SwinUNETR architectures are taken directly from MONAI's model
zoo and are \emph{not} our original designs.
No pre-trained weights were used; all models were trained from scratch.

\paragraph{Dataset.}
The Medical Segmentation Decathlon Task~01 dataset \cite{antonelli2022}
is a publicly available benchmark used without modification.

\paragraph{Original contributions.}
The following are our original work:
\begin{itemize}[leftmargin=*]
    \item The \texttt{BoundaryAwareLoss} module combining DiceCE with
          periodic HausdorffDT (every $k=10$ steps), implemented in
          \texttt{src/losses.py}.
    \item The frequency-throttled training loop modification in
          \texttt{src/train.py}.
    \item The controlled experimental design, data pipeline
          (supporting both MSD and BraTS2021 formats), evaluation harness,
          and all reported experiments.
    \item All figures and analyses presented in this report.
\end{itemize}

All code is available at:
\url{https://github.com/devvansh1997/3D-BrainTumor-Seg}

% ─────────────────────────────────────────────────────────────────────────────
\bibliographystyle{plainnat}
\bibliography{references}

\end{document}
